\section{量角器}
基于 Yves Combe 的思路，以下宏用于绘制量角器。
原理很简单：只需指定一条射线。量角器放置在其原点 $O$ 处，射线的方向由 $A$ 给出。角度按逆时针方向测量；使用选项 \tkzname{return} 时，量角器被反转，角度按顺时针方向测量。
\subsection{宏 \tkzcname{tkzProtractor}} % (fold)
\label{sub:the_macro_tkzcname_tkzprotractor}

% subsection the_macro_tkzcname_tkzprotractor (end)
\begin{NewMacroBox}{tkzProtractor}{\oarg{local options}\parg{$O,A$}}%
\begin{tabular}{lll}%
选项    & 默认值 & 说明     \\ 
\midrule
\TOline{lw}  {0.4 pt} {线宽}
\TOline{scale}  {1} {比例：调整量角器的大小} 
\TOline{return} {false} {三角函数圆反向}
\end{tabular}
\end{NewMacroBox}

\subsubsection{圆形量角器} 
逆时针测量：

\begin{tkzexample}[latex=7cm,small] 
\begin{tikzpicture}[scale=.5]
\tkzDefPoint(2,0){A}\tkzDefPoint(0,0){O}
\tkzDefShiftPoint[A](31:5){B}
\tkzDefShiftPoint[A](158:5){C}
\tkzDrawPoints(A,B,C)
\tkzDrawSegments[color = red,
    line width = 1pt](A,B A,C)
  \tkzProtractor[scale = 1](A,B) 
\end{tikzpicture}
\end{tkzexample}  

\subsubsection{圆形量角器，透明且反向}

\begin{tkzexample}[latex=7cm,small] 
\begin{tikzpicture}[scale=.5]
  \tkzDefPoint(2,3){A}
  \tkzDefShiftPoint[A](31:5){B}  
   \tkzDefShiftPoint[A](158:5){C}   
  \tkzDrawSegments[color=red,line width=1pt](A,B A,C)
  \tkzProtractor[return](A,C) 
\end{tikzpicture}
\end{tkzexample}
\endinput
